\documentclass{article}

\usepackage[italian]{babel}
\usepackage[margin=20mm, footskip = 20pt]{geometry}
\usepackage{graphicx}
\usepackage{subfiles}
\usepackage{hyperref}
\usepackage{nameref}
\usepackage{titlesec}
\usepackage{longtable}
\usepackage[table]{xcolor}
\usepackage{titling}
\usepackage{lastpage}
\usepackage{ifthen}
\usepackage{calc}
\usepackage{soulutf8}
\usepackage{contour}
\usepackage{float}
\usepackage{fancyhdr}
\usepackage{multirow}
\usepackage{pgfgantt}
\usepackage{lscape}

% definizione dei percorsi in cui cercare immagini
\graphicspath{ {./}
    {./img/}
}

% setup della sottolineatura
\setuldepth{Flat}
\contourlength{0.8pt}

\newcommand{\uline}[1]{%
  \ul{{\phantom{#1}}}%
  \llap{\contour{white}{#1}}%
}

% setup dei link
\hypersetup{
  % set true if you want colored links (instead of boxes)
  colorlinks=true,
  % set to all if you want both sections and subsections linked
  linktoc=all,
  % set color for file links
  filecolor=blue,
  % set color for internal links
  linkcolor=black,
  % set url color
  urlcolor=blue,
  % set characters encoding in the bookmarks tab
  pdfencoding=unicode,
}

% setup forma \paragraph e \subparagraph
\titleformat{\paragraph}[hang]{\normalfont\normalsize\bfseries}{\theparagraph}{1em}{}
\titleformat{\subparagraph}[hang]{\normalfont\normalsize\bfseries}{\thesubparagraph}{1em}{}

% setup profondità indice di default
\setcounter{secnumdepth}{5}
\setcounter{tocdepth}{5}

\makeatletter %% non togliere, i comandi che definiscono i placeholder vanno qui
% esempio di utilizzo: \appendToGraphicspath{./img/} (un comando diverso per ogni path da includere)
% N.B.: ci DEVE essere un forward slash alla fine del path, a indicare che è una cartella.
\newcommand\appendToGraphicspath[1]{%
  \g@addto@macro\Ginput@path{{#1}}%
}

\newcommand{\setTitle}[1]{%
  \newcommand{\@placeholderTitle}{#1}%
}
\newcommand{\placeholderTitle}{\@placeholderTitle}

\newcommand{\setUso}[1]{%
  \newcommand{\@uso}{#1}%
}
\newcommand{\uso}{\@uso}

\newcommand{\setVersione}[1]{%
  \newcommand{\@versione}{#1}%
}
\newcommand{\versione}{\@versione}

\newcommand{\disabilitaVersione}{%
  \renewcommand{\setVersione}[1]{}%
  \renewcommand{\versione}{DISABILITATA}
}

\newcommand{\setResponsabile}[1]{%
  \newcommand{\@responsabile}{#1}%
}
\newcommand{\responsabile}{\@responsabile}

\newcommand{\setRedattori}[1]{%
  \newcommand{\@redattori}{#1}%
}
\newcommand{\redattori}{\@redattori}

\newcommand{\setVerificatori}[1]{%
  \newcommand{\@verificatori}{#1}%
}
\newcommand{\verificatori}{\@verificatori}

\newcommand{\setDescrizione}[1]{%
  \newcommand{\@descrizione}{#1}%
}
\newcommand{\descrizione}{\@descrizione}

\newcommand{\setModifiche}[1]{%
  \newcommand{\@modifiche}{#1}%
}

\newcommand{\modifiche}{\@modifiche}
\makeatother %% non togliere, i comandi che definiscono i placeholder vanno qui

% hook per lo script che genera il glossario
\newcommand{\glossario}[1]{\underline{#1}\textsubscript{G}}

% comandi per rendere opzionali gli elenchi di figure
\newcommand{\elencoFigure}{%
  \renewcommand{\listfigurename}{Elenco delle figure}%
  \listoffigures%
}

\newcommand{\disabilitaElencoFigure}{%
  \renewcommand{\elencoFigure}{}%
}

% comandi per rendere opzionali le tabelle
\newcommand{\elencoTabelle}{%
  \renewcommand{\listtablename}{Elenco delle tabelle}%
  \listoftables%
}

\newcommand{\disabilitaElencoTabelle}{%
  \renewcommand{\elencoTabelle}{}%
}

%Qui ci andrà il percorso delle immagini da includere in analisi dei requisiti

\appendToGraphicspath{../../commons/img/}

%Tutti questi set permettono di modificare in modo adatto i placeholder nel template
\setTitle{Glossario}

\setVersione{v1.0.0}

\setResponsabile{Ian Nicolas Di Menna}

\setRedattori{Sofia Bononi \\ &
Enrico Buratto \\ &
Enrico Galdeman \\ &
Nicholas Miazzo \\ &
Marco Nardellotto}

\setVerificatori{Enrico Buratto \\ &
Marco Nardellotto}

\setUso{Esterno}

\setDescrizione{Documento con lo scopo di chiarire le possibili ambiguità tra i termini utilizzati all'interno dei documenti scritti e redatti dal team.}

\disabilitaElencoFigure{}
\disabilitaElencoTabelle{}

\setModifiche{v1.0.0 & Ian Nicolas Di Menna & Responsabile & 2020-04-12 & Approvazione del documento \\
v0.2.0 & Enrico Buratto & Verificatore & 2020-04-11 & Verifica termini \textsc{Piano di Qualifica} e \textsc{Norme di Progetto} \\
v0.1.0 & Marco Nardellotto & Verificatore & 2020-04-10 & Verifica termini \textsc{Studio di Fattibilità},\textsc{Piano di Progetto} e \textsc{Analisi dei requisiti} \\
v0.0.6 & Enrico Galdeman & Analista & 2020-04-09 & Aggiunta termini \textsc{Norme di Progetto} \\
v0.0.5 & Nicholas Miazzo & Analista & 2020-04-09 & Aggiunta termini \textsc{Analisi dei Requisiti} \\
v0.0.4 & Enrico Buratto & Analista & 2020-04-08 & Aggiunta termini \textsc{Piano di Progetto} \\
v0.0.3 & Sofia Bononi & Analista & 2020-03-31 & Aggiunta termini \textsc{Piano di Qualifica} \\
v0.0.2 & Marco Nardellotto & Analista & 2020-03-30 & Aggiunta termini \textsc{Studio di Fattibilità} \\
v0.0.1 & Alessandro Franchin & Responsabile & 2020-03-16 & Creazione del documento}

\begin{document}

\pagenumbering{gobble}

\begin{titlepage}% per non stampare il numero della pagina

  \raggedleft% allinea a destra la pagina
  \rule{1pt}{\textheight}% linea verticale
  \hspace{0.05\textwidth}% spazio tra linea e testo
  % lasciare questa riga per il corretto funziomento di \parbox
  \parbox[b]{0.75\textwidth}{% paragrafo che tiene il testo a destra della riga cambiando la larghezza il testo si muove a destra o a sinistra
  {\hspace{0.15\textwidth}\includegraphics[width=4cm,height=4cm]{logo.jpg}}\\[4\baselineskip] % logo
  {\Huge\bfseries CoffeeCode \\[1\baselineskip] Predire in Grafana}\\[6\baselineskip] % titolo
  {\Large\textsc{\placeholderTitle{}}}\\[7\baselineskip] % nome del documento
  {\begin{tabular}{r l}
    % testo in grassetto
    \textbf{Versione}     & \versione{}               \\
    \textbf{Approvazione} & \responsabile{}           \\
    \textbf{Redazione}    & \redattori{}              \\
    \textbf{Verifica}     & \verificatori{}           \\
    \textbf{Uso}          & \uso{}                    \\
    \textbf{Destinato a}  & CoffeeCode                \\
                          & prof.\ Vardanega Tullio   \\
                          & prof.\ Cardin Riccardo    \\
    \ifthenelse{\equal{\uso}{Esterno}}{
                          & Zucchetti S.p.A.       \\
    }{}
  \end{tabular}}\\[5\baselineskip]

  {\bfseries Descrizione}\\
  {\descrizione{}}\\[2\baselineskip]
  {\texttt{coffeecodeswe@gmail.com}}\\[\baselineskip] % email
  }

\end{titlepage}

\newgeometry{textheight=660pt, lmargin=2cm, tmargin=2cm, rmargin=2cm}

% setup di header e footer nelle pagine senza numero
\fancypagestyle{nopage}{%
  \fancyhf{}%
  \fancyhead[L]{\includegraphics[width=1.3cm]{logo.jpg}}%
  \fancyhead[R]{\emph{CoffeeCode}\\\placeholderTitle{}}%
}
% setup di header e footer nelle pagine col numero
\fancypagestyle{usual}{%
  \fancyhf{}%
  \fancyhead[L]{\includegraphics[width=1.3cm]{logo.jpg}}%
  \fancyhead[R]{\emph{CoffeeCode}\\\placeholderTitle{}}%
  \fancyfoot[R]{\thepage\ di~\pageref{LastPage}}%
}
\setlength{\headheight}{1.8cm}

\newpage
\pagestyle{nopage}

\section*{Registro delle modifiche}%
\label{sec:registro_delle_modifiche}

\rowcolors{2}{white!80!lightgray!90}{white}
\renewcommand{\arraystretch}{2} % allarga le righe con dello spazio sotto e sopra
\begin{longtable}[H]{>{\centering\bfseries}m{2cm} >{\centering}m{5cm} >{\centering}m{2.5cm} >{\centering\arraybackslash}m{6.5cm}}
  \caption{Registro delle modifiche}%
  \label{tab:registro_delle_modifiche}                                                    \\
  \rowcolor{lightgray}
  {\textbf{Versione}} & {\textbf{Nominativo}} & {\textbf{Data}} & {\textbf{Descrizione}}  \\
  \endfirsthead%
  \rowcolor{lightgray}
  {\textbf{Versione}} & {\textbf{Nominativo}} & {\textbf{Data}} & {\textbf{Descrizione}}  \\
  \endhead%
  \modifiche{}%
\end{longtable}
% section registro_delle_modifiche (end)

\newpage
\thispagestyle{nopage}
\pagenumbering{roman}
\tableofcontents

\elencoFigure{}%

\elencoTabelle{}%

\newpage

\pagestyle{usual}
\pagenumbering{arabic}


\section{A}
\subfile{components/A.tex}
\newpage

\section{B}
\subfile{components/B.tex}
\newpage

\section{C}
\subfile{components/C.tex}
\newpage

\section{D}
\subfile{components/D.tex}
\newpage

\section{E}
\subfile{components/E.tex}
\newpage

\section{F}
\subfile{components/F.tex}
\newpage

\section{G}
\subfile{components/G.tex}
\newpage

\section{H}
\subfile{components/H.tex}
\newpage

\section{I}
\subfile{components/I.tex}
\newpage

\section{J}
\subfile{components/J.tex}
\newpage

\section{K}
\subfile{components/K.tex}
\newpage

\section{L}
\subfile{components/L.tex}
\newpage

\section{M}
\subfile{components/M.tex}
\newpage

\section{N}
\subfile{components/N.tex}
\newpage

\section{O}
\subfile{components/O.tex}
\newpage

\section{P}
\subfile{components/P.tex}
\newpage

\section{Q}
\subfile{components/Q.tex}
\newpage

\section{R}
\subfile{components/R.tex}
\newpage

\section{S}
\subfile{components/S.tex}
\newpage

\section{T}
\subfile{components/T.tex}
\newpage

\section{U}
\subfile{components/U.tex}
\newpage

\section{V}
\subfile{components/V.tex}
\newpage

\section{W}
\subfile{components/W.tex}
\newpage

\section{X}
\subfile{components/X.tex}
\newpage

\section{Y}
\subfile{components/Y.tex}
\newpage

\section{Z}
\subfile{components/Z.tex}

\end{document}
